\documentclass[12pt, a4paper]{article}
\usepackage[utf8]{inputenc}
\usepackage{geometry}
\geometry{a4paper, margin=1in}
\usepackage{hyperref}
\usepackage{graphicx}
\usepackage{listings}
\usepackage{xcolor}

\title{\textbf{ACSAR-E: Real-World Implementation \& Simulation Mapping Proposal}}
\author{Autonomous Unitree Go2 Security \& Patrol Project}
\date{\today}

\begin{document}

\maketitle

\section*{Overview}
\textbf{Phase:} Simulation Validation \& Hardware Translation

\section{Spatial Awareness \& Conditional Zone Detection}
In our \textit{Gazebo-Lite} simulation, we utilized mathematical 3D bounding boxes (AABB) and Three.js Raycasting to determine if the robot entered ``Restricted Zones.'' To deploy this in real life:

\begin{itemize}
    \item \textbf{Real-World Implementation (Geofencing via SLAM):} We will use the Unitree Go2's onboard \textbf{SLAM (Simultaneous Localization and Mapping)} capabilities. The robot generates a persistent 3D point-cloud map of the physical office.
    \item \textbf{Coordinate Translation:} We map our physical office blueprint to an absolute coordinate system (e.g., origin $(0,0)$ at the charging dock). Conditional zones (like the \textit{High Voltage Lab}) are defined as Cartesian boundary polygons $[(x_1, z_1), (x_2, z_2)...]$ in a Python config file.
    \item \textbf{Trigger Logic:} The real robot's odometry data continuously feeds its current $[X, Y, Z]$ position to our Python script. A simple Point-in-Polygon (PiP) algorithm checks if the robot's real-time coordinate falls inside a restricted polygon, instantly triggering the policy violation alerts.
\end{itemize}

\section{Dynamic Obstacle Avoidance \& Human Interaction (YOLO)}
Our current simulation uses A* to navigate static walls, and a LiDAR failsafe to stop before crashing. In a real office, humans move dynamically.

\begin{itemize}
    \item \textbf{The Problem:} A* calculates paths assuming the world is frozen. If a human walks into the corridor, A* will fail or the robot will abruptly emergency-stop.
    \item \textbf{The Solution (YOLOv8 + LiDAR Fusion):} We will intercept the RTSP video stream from the Go2's front-facing camera and feed it into a \textbf{YOLOv8 Object Detection} neural network running on an edge-compute node (like a Jetson Orin or the onboard high-performance computer).
    \item \textbf{How it Works:}
    \begin{enumerate}
        \item YOLO detects a bounding box labeled \texttt{person}.
        \item We overlay that bounding box onto the hardware LiDAR depth-map to extract the precise distance and vector of the human.
        \item We dynamically inject a high-cost ``temporary obstacle blob'' into the robot's local Costmap.
        \item The robot smoothly arcs around the moving human instead of stopping, utilizing a local trajectory planner (like DWA - Dynamic Window Approach) while maintaining its global A* trajectory.
    \end{enumerate}
\end{itemize}

\section{Advanced Routing: Water Flow vs. Quantum-Inspired Search}
While A* is perfect for Point A to Point B, patrolling a dynamic security environment requires more organic algorithms.

\begin{itemize}
    \item \textbf{Water Flow Algorithm (Wavefront/Brushfire):} Instead of calculating a single thin line, the Water Flow algorithm simulates pouring water from the robot's location. The ``water'' fills the corridors based on terrain costs. This is exceptionally powerful for \textbf{Exploration and Sweeping}. If we need the robot to search an entire room for anomalies, Water Flow guarantees 100\% floor-plan coverage without redundant overlapping.
    \item \textbf{Quantum-Inspired Search Optimization:} If the robot needs to patrol 50 different high-risk waypoints (Server racks, doors, windows), finding the absolute fastest route is the ``Traveling Salesman Problem'' (highly computationally expensive). Quantum-inspired algorithms (like Q-Learning or Quantum Annealing simulations) evaluate multiple routing probabilities simultaneously. We can deploy a cloud-based quantum-inspired solver to recalculate the robot's patrol route every 5 minutes based on dynamic heatmaps (e.g., routing the robot to areas where YOLO detected the most human foot traffic earlier in the day).
\end{itemize}

\section{Phase Execution: From Simulation to Hardware}
Our Python simulation architecture was intentionally designed to perfectly mirror the real-world deployment.

\subsection{Phase 1: Virtual Validation (Complete)}
We have successfully built \texttt{unitree\_sim.py}, which acts as a virtual SDK. Our \texttt{test\_run.py} script proves that a standalone Python loop can ingest real-time LiDAR telemetry, calculate distances, and send asynchronous movement commands to a 3D robot over WebSockets.

\subsection{Phase 2: Hardware API Bridge (Upcoming)}
To run \texttt{test\_run.py} on the physical Unitree Go2, we only need to swap the \texttt{Go2Sim} import with the official Unitree SDK.

\begin{lstlisting}[language=Python, basicstyle=\ttfamily\small, frame=single, breaklines=true]
# Simulation Code
from unitree_sim import Go2Sim 
robot = Go2Sim()

# Real-Life Code (Swapped)
from unitree_sdk2 import Go2Hardware
robot = Go2Hardware(ip="192.168.123.161")
\end{lstlisting}

Because we adhered to standard generic function calls (\texttt{get\_lidar\_front()}, \texttt{get\_position()}, \texttt{move\_to()}), the Python logic \textbf{does not need to change}. The official Unitree SDK will simply translate our \texttt{move\_to(x, z)} commands into native CAN-bus motor torques.

\subsection{Phase 3: Edge AI Deployment}
\begin{enumerate}
    \item Load the mapped office SLAM data into the Go2.
    \item Deploy the \texttt{test\_run.py} (with YOLO integrated) onto an NVIDIA Jetson mounted to the robot's back.
    \item The Jetson handles heavy YOLO processing and Water Flow routing, feeding lightweight movement vectors directly into the Go2's movement controller.
\end{enumerate}

\end{document}
